% GENERATED FILE. Edit canonical lessons, then run npm run generate:books.
% canonical-source-hash: fnv1a64:aa5c7e544198efe0
% canonical-lessons: PT-C78-chapeu, PT-C78-casaco, PT-C78-saia, PT-C78-meias, PT-C78-cinto

\chapter{Hat, Coat, Skirt, Socks, Belt}
\label{ch:pt-hat-coat-skirt-socks-belt}
\hlchaptermodality{78}

\begin{chapteropening}
\textbf{By the end of this chapter:} \emph{I can name a hat, a coat, a skirt, socks and a belt.}

Complete the last lesson of chapter 78: I can name a hat, a coat, a skirt, socks and a belt.
\end{chapteropening}

\section[o chapéu]{o chapéu — a hat}

\label{lesson:PT-C78-chapeu}



\begin{quote}
Before the new one: say the Portuguese for trousers, then the Portuguese for shoes.
\end{quote}

\subsection*{You'll want to know: o chapéu}
\textbf{o chapéu} — \textquotedblleft{}a hat\textquotedblright{}.

\begin{cousinweb}
\textbf{chapéu} $\leftarrow$ Old French \emph{chapel}, \textquotedblleft{}hat\textquotedblright{} — borrowed. Behind it is Vulgar Latin \emph{*cappellus}, from Late Latin \emph{cappa}, \textquotedblleft{}hooded cloak\textquotedblright{} (EtymDB 2.1).
\end{cousinweb}

\begin{grammarlens}[title={What you've built}]
One more word for this chapter.
\end{grammarlens}

\subsection*{Your turn}
\begin{itemize}
\raggedright
  \item \emph{Say it:} \emph{o chapéu}
  \item \emph{Say it:} \emph{o chapéu}, once more
  \item \emph{Say it:} say \emph{o chapéu}
  \item \emph{Recall:} say the Portuguese for trousers, then the Portuguese for shoes, then say \emph{o chapéu} again
\end{itemize}

\begin{tcolorbox}[breakable,colback=teal!4,colframe=teal!35!black,title={Before you move on}]
What does o chapéu mean? (\textquotedblleft{}A hat\textquotedblright{}.) Say it once more.
\end{tcolorbox}
\section[o casaco]{o casaco — a coat}

\label{lesson:PT-C78-casaco}



\begin{quote}
Before the new one: say the Portuguese for shoes, then the Portuguese for a hat.
\end{quote}

\subsection*{You'll want to know: o casaco}
\textbf{o casaco} — \textquotedblleft{}a coat\textquotedblright{}.

\begin{cousinweb}
\textbf{casaco} is a masculine of an older word for a coat (EtymDB 2.1), borrowed from French \emph{casaque} (EtymDB 2.1).
\end{cousinweb}

\begin{grammarlens}[title={What you've built}]
One more word for this chapter.
\end{grammarlens}

\subsection*{Your turn}
\begin{itemize}
\raggedright
  \item \emph{Say it:} \emph{o casaco}
  \item \emph{Say it:} \emph{o casaco}, once more
  \item \emph{Say it:} say \emph{o casaco}
  \item \emph{Recall:} say the Portuguese for shoes, then the Portuguese for a hat, then say \emph{o casaco} again
\end{itemize}

\begin{tcolorbox}[breakable,colback=teal!4,colframe=teal!35!black,title={Before you move on}]
What does o casaco mean? (\textquotedblleft{}A coat\textquotedblright{}.) Say it once more.
\end{tcolorbox}
\section[a saia]{a saia — a skirt}

\label{lesson:PT-C78-saia}



\begin{quote}
Before the new one: say the Portuguese for a hat, then the Portuguese for a coat.
\end{quote}

\subsection*{You'll want to know: a saia}
\textbf{a saia} — \textquotedblleft{}a skirt\textquotedblright{}.

\begin{cousinweb}
\textbf{saia} $\leftarrow$ Vulgar Latin \emph{*sagia} (EtymDB 2.1), from Latin \emph{sagum}, \textquotedblleft{}a soldier's cloak\textquotedblright{} — inherited.
\end{cousinweb}

\begin{grammarlens}[title={What you've built}]
One more word for this chapter.
\end{grammarlens}

\subsection*{Your turn}
\begin{itemize}
\raggedright
  \item \emph{Say it:} \emph{a saia}
  \item \emph{Say it:} \emph{a saia}, once more
  \item \emph{Say it:} say \emph{a saia}
  \item \emph{Recall:} say the Portuguese for a hat, then the Portuguese for a coat, then say \emph{a saia} again
\end{itemize}

\begin{tcolorbox}[breakable,colback=teal!4,colframe=teal!35!black,title={Before you move on}]
What does a saia mean? (\textquotedblleft{}A skirt\textquotedblright{}.) Say it once more.
\end{tcolorbox}
\section[as meias]{as meias — socks}

\label{lesson:PT-C78-meias}



\begin{quote}
Before the new one: say the Portuguese for a coat, then the Portuguese for a skirt.
\end{quote}

\subsection*{You'll want to know: as meias}
\textbf{as meias} — \textquotedblleft{}socks\textquotedblright{}.

\begin{cousinweb}
\textbf{meias} is short for \textquotedblleft{}half-hose\textquotedblright{}: the word is the feminine of the word for \textquotedblleft{}half\textquotedblright{}, which $\leftarrow$ Latin \emph{medius}, \textquotedblleft{}middle\textquotedblright{} (Lewis \& Short).
\end{cousinweb}

\begin{grammarlens}[title={What you've built}]
One more word for this chapter.
\end{grammarlens}

\subsection*{Your turn}
\begin{itemize}
\raggedright
  \item \emph{Say it:} \emph{as meias}
  \item \emph{Say it:} \emph{as meias}, once more
  \item \emph{Say it:} say \emph{as meias}
  \item \emph{Recall:} say the Portuguese for a coat, then the Portuguese for a skirt, then say \emph{as meias} again
\end{itemize}

\begin{tcolorbox}[breakable,colback=teal!4,colframe=teal!35!black,title={Before you move on}]
What does as meias mean? (\textquotedblleft{}Socks\textquotedblright{}.) Say it once more.
\end{tcolorbox}
\section[o cinto]{o cinto — a belt}

\label{lesson:PT-C78-cinto}



\begin{quote}
Before the new one: say the Portuguese for a skirt, then the Portuguese for socks.
\end{quote}

\subsection*{You'll want to know: o cinto}
\textbf{o cinto} — \textquotedblleft{}a belt\textquotedblright{}.

\begin{cousinweb}
\textbf{cinto} $\leftarrow$ Latin \emph{cinctus}, \textquotedblleft{}girdle, belt\textquotedblright{} (EtymDB 2.1), from \emph{cingere}, \textquotedblleft{}to gird\textquotedblright{} — inherited.
\end{cousinweb}

\begin{grammarlens}[title={What you've built}]
One more word for this chapter.
\end{grammarlens}

\subsection*{Your turn}
\begin{itemize}
\raggedright
  \item \emph{Say it:} \emph{o cinto}
  \item \emph{Say it:} \emph{o cinto}, once more
  \item \emph{Say it:} say \emph{o cinto}
  \item \emph{Recall:} say the Portuguese for a skirt, then the Portuguese for socks, then say \emph{o cinto} again
\end{itemize}

\begin{tcolorbox}[breakable,colback=teal!4,colframe=teal!35!black,title={Before you move on}]
What does o cinto mean? (\textquotedblleft{}A belt\textquotedblright{}.) Say it once more.
\end{tcolorbox}
